\chapter{Типовые игровые данные, подлежащие хранению, обновлению и передаче}

В процессе работы многопользовательского игрового приложения сервер и клиенты обмениваются различными
типами данных\cite{gregory2018gamedev}, описывающих состояние объектов, игровые события и вспомогательные параметры. Эти данные
определяют внутреннюю логику игрового мира и корректность взаимодействия игроков.

\section{Данные состояния объектов}

Состояние объектов, которое описывает характеристики сущностей, является основой игрового мира. Эти данные подлежат частому обновлению.

К типичным данным относятся:
\begin{itemize}
    \item позиция объекта;
    \item скорость и направление движения;
    \item повороты;
    \item состояние анимации;
    \item физические параметры.
\end{itemize}

Эти данные требуют высокой точности и частой синхронизации. Особенно это критично для жанров FPS, гонок и экшен-игр,
где любое отклонение в позиции приводит к заметным ошибкам восприятия.

\section{Данные игровой логики}

Помимо визуального состояния объектов существуют параметры, влияющие на игровую логику.

В эту категорию входят:
\begin{itemize}
    \item здоровье;
    \item ресурсы;
    \item эффекты;
    \item флаги состояний;
\end{itemize}

Данные игровой логики обновляются реже, чем движение, но требуют надёжной доставки, поскольку ошибки в этих данных
приводят к различию в результатах игровых событий.

\section{Данные взаимодействия и событий}

Игровые события отображают дискретные изменения состояния мира. Они могут быть результатом действий игрока или
следствием внутренних механик.

Примеры событий:
\begin{itemize}
    \item совершение действия;
    \item взаимодействие с объектом;
    \item создание или уничтожение объекта;
    \item триггеры.
\end{itemize}

События требуют надёжной доставки, поскольку их потеря приводит к
расхождениям в симуляции между игроками.

\section{Вспомогательные данные}

Некоторые игры используют дополнительные данные, не относящиеся напрямую к состоянию объектов, но влияющие на
корректность симуляции.

К таким данным относятся:
\begin{itemize}
    \item навигационные сетки;
    \item вспомогательные маркеры;
    \item данные о временных параметрах.
\end{itemize}

Эти данные чаще хранятся на сервере и передаются клиенту при необходимости.

\section{Визуальные данные}

Некоторые данные не влияют на игровую логику напрямую и могут синхронизироваться реже или даже генерироваться на
клиенте без участия сервера.

К таким данным относятся:
\begin{itemize}
    \item эффекты частиц;
    \item косметические параметры;
    \item звуковые события.
\end{itemize}

Передача таких данных оптимизируется в первую очередь при ограничении пропускной способности.

\section{Метаданные и технические данные}

Помимо игровых данных существует набор технической информации, которая обеспечивает корректную работу сетевого
взаимодействия.

К ним относятся:
\begin{itemize}
    \item идентификаторы объектов;
    \item временные метки;
    \item контрольные суммы;
    \item служебные флаги.
\end{itemize}

Такие данные не видны игроку, но они крайне важны для поддержания согласованности состояния и предотвращения ошибок.

